\documentclass[12pt]{article}
\usepackage{graphics,graphicx,array,url}

\headheight=0in
\headsep=0.2in
\topskip=0in
\setlength{\oddsidemargin}{-0.5in}
\setlength{\evensidemargin}{0cm}
\setlength{\topmargin}{-0.5in}
\setlength{\textwidth}{7.0in}
\setlength{\textheight}{9in}

\begin{document}
\hfill \includegraphics[width=0.35\textwidth]{siena_phys_astro_print_crop.jpg}

\vspace{0.2cm}
\begin{center}
{\LARGE {\bf Frequently Asked Questions}} \\
\medskip
{\Large {\bf 4/1 Program with Clarkson University}} 
\end{center}
\vspace{0.3cm}

%\noindent This is good.

\begin{itemize}
\item{{\bf {\em What is this document? Where can I get for more information?}} 

This document is meant to be an unofficial guide for students at Siena College
\emph{considering} the 4/1 Program; it is not a replacement for speaking at
length with your academic advisor and with the Dean of the School of Engineering
and Computer Science at Clarkson University, Robert Kozik (bkozik@clarkson.edu).
At Siena you may also reach out to Prof. John Moustakas (jmoustakas@siena.edu)
in the Physics Department if you have specific questions not answered below.
Finally, you should reach out to your peers at Siena who have already started
taking classes at Clarkson as part of the 4/1 Program---their experiences are an
invaluable resource!}

\item{{\bf {\em What is the 4/1 program?}}

The 4/1 Program is an articulation between Siena College and Clarkson University
which allows undergraduate students pursuing their BS in Physics, Environmental
Science, Computer Science, or Mathematics to take graduate-level courses at the
Clarkson University Capital Region Campus in Schenectady, NY, giving them a jump
on obtaining their engineering masters degree in just one additional year
(hence, the ``4/1 Program'').  You may take up to \underline{three} courses at
Clarkson as part of your Siena tuition (i.e., without paying any additional
fees).  In addition, the courses you take at Clarkson may also satisfy your
upper-level major course requirements, although that depends on the course and
your department---see your advisor for details.}

%In the physics department these courses also count toward your
%required upper-level physics courses (as PHYS400).  See your advisor for
%details.} 

\item{{\bf {\em What engineering degrees are available?}}

The three Master of Science (M.S.) programs currently available to students are
{\em Electrical Engineering} (MSEE), {\em Engineering and Management Systems}
(MSEM), and {\em Energy Systems} (MSE).  Quoting from the program catalog, the
MSEE program ``explores technologies and industry opportunities in modern
electric machinery, modeling and control of power electronics.''  The MSEM
program integrates ``engineering and/or information systems technologies with
the core components of an MBA, including Business of Energy Program courses.''
And the MSE degree enables ``students to integrate (1) mechanical/electrical
energy related courses, (2) mechanical and electrical fundamental discipline
courses, and (3) non-technical courses regarding the impact of environmental,
economic, and regulatory issues on energy.''}

\item{{\bf {\em When can I start taking classes at Clarkson?}}

Students can start taking classes as soon as they have room in their schedule
and have sufficient mathematics preparation, typically in their junior or senior
year.  Although these are not formal prerequisites, for the more technical
electrical engineering courses (e.g., {\em Linear Control Systems}) it is
strongly recommended that students have completed all three semesters of
calculus, differential equations, and either linear algebra or applied
mathematics (ideally, both), as well as some experience with {\tt Matlab}/{\tt
  Simulink}.  Some of the business-oriented courses (e.g., {\em Fundamentals of
  the Business of Energy}) may be taken without having completed as much math,
although they are still graduate-level courses and require a significant amount
of work.}

\item{{\bf {\em When are the academic terms at Clarkson?}}

Unlike Siena, Clarkson is on a 10-week trimester system.  The fall, winter, and
spring trimesters are roughly from early September to late November, early
January to mid-March, and early April to early June, respectively.  (Please note
that only the fall term is aligned with the housing term at Siena; therefore, if
you do not live locally and need to take winter and/or spring courses you need
to consider where you will live while taking courses at Clarkson.)  In addition,
there is typically a summer session of selective classes from mid-June to the
end of August.  The classes typically take place during the week and in the
evening.  The average class size is around 12 students.}

\item{{\bf {\em What classes are offered?  What are the prerequisites?}}

The courses available varies somewhat from year-to-year, so you should speak
with your advisor and Dean Kozik about which courses are available and whether
you are ready to start taking courses at Clarkson.  Some of the courses our
students have taken (and when they are typically offered) include {\em Linear
  Control Systems} (fall), {\em Fundamentals of Business of Energy} (online,
offered frequently), {\em Solar Energy Engineering} (winter), {\em Fuel Cell \&
  Hydrogen Engineering} (varies), {\em Wind Energy Engineering} (project-based,
spring), and {\em Electromechanical Energy Conversion}.}

\item{{\bf {\em How do I apply to the program?  How do I register for classes?
      What are the minimum requirements?}}

In order to apply to the program you must submit an \underline{Application
  Form}\footnote{http://www.clarkson.edu/artsandsci/documents/graddocuments/CompleteApplicationPacket.pdf},
an official transcript, a personal statement, proof of immunizations, and two
(not three) letters of recommendation; the application fee is waived for Siena
students.  Note that a minimum 3.0 GPA is also strongly recommended.  It is
possible---and indeed recommended---that you take at least one class
\emph{before} officially applying to Clarkson, so you can get a better idea of
the courses and the program itself.  In that case you should get in contact with
Dean Kozik about what documentation must be received in order to take an initial
class as a non-matriculated student.

In order to register for classes you need to fill out the
\underline{Registration
  Form}\footnote{\url{http://clarkson.edu/sas/forms/crc_registration_siena.pdf}}
specifically designed for Siena students.  The second page of this registration
form has detailed instructions that you should follow carefully, but basically
the form will need to be signed by your acadmemic advisor, the Dean of the
School of Science, and the Siena College registrar (Jim Serbalik, Siena Hall
102) before it is sent to Clarkson.  Be sure to get these signatures well in
advance of the start of classes!}

\item{{\bf {\em How do I get credit for these classes on my CAPP report?}}

The classes you take at Clarkson as part of the 4/1 Program are treated by the
Siena registrar as if they were courses you took at Siena itself.  Therefore,
once you complete the course at Clarkson and your grade has been submitted, a
transcript from Clarkson will be sent to Siena so the grade and course can
appear on your transcript and CAPP report.  This is supposed to happen
automatically between the two institutions (without a fee!); however, it
behooves you to follow up in person, by calling, or by emailing---and with your
advisor---to make sure all the necessary information has been received and that
your CAPP report looks correct.  For physics students, courses taken at Clarkson
appear on their transcript and CAPP report as PHYS400 ({\em Special Topics in
  Physics}).}

\item{{\bf {\em Are there research opportunities?}}

Unfortunately there are no research opportunities, although you should speak
with Dean Kozik early and often about possible paid and unpaid internships.}

%\item{{\bf {\em How do I get to Clarkson and where can I park?}}
%Answer}

%\item{{\bf {\em Where can I read the articulation agreement between Siena and
%      Clarkson?}}
%link}

\item{{\bf {\em What happened to Union Graduate College?}}

In January 2016 Union Graduate College became the Clarkson University Capital
Region Campus, a satellite campus of Clarkson University which is located in
Potsdam, NY.}
  
\end{itemize}
  
\end{document}
